\documentclass[12pt]{article}
\usepackage[T1]{fontenc}
\usepackage[utf8]{inputenc}
\usepackage{lmodern}
\usepackage{textcomp}
\usepackage{enumitem}
\usepackage{geometry}
\usepackage{graphicx}
\usepackage{amsmath, amssymb}
\geometry{margin=1in}
\begin{document}
\begin{center}
Physics - Undergraduate Level Assessment 20 \
Total Marks: 40 \
Answer all questions.
\end{center}

\section*{Multiple Choice (10 marks)}
\begin{enumerate}[label=\arabic*.]
\item Which of the following lasers utilizes transitions that involve the energy levels of free atoms?
\begin{enumerate}[label=(\alph*)]
    \item Diode laser
    \item Dye laser
    \item Free-electron laser
    \item Gas laser
\end{enumerate}
\item A refracting telescope consists of two converging lenses separated by 100 cm. The eye- piece lens has a focal length of 20 cm. The angular magnification of the telescope is
\begin{enumerate}[label=(\alph*)]
    \item 4
    \item 5
    \item 6
    \item 20
\end{enumerate}
\item A resistor in a circuit dissipates energy at a rate of 1 W. If the voltage across the resistor is doubled, what will be the new rate of energy dissipation?
\begin{enumerate}[label=(\alph*)]
    \item 0.25 W
    \item 0.5 W
    \item 1 W
    \item 4 W
\end{enumerate}
\item A single-electron atom has the electron in the l = 2 state. The number of allowed values of the quantum number m\_l is
\begin{enumerate}[label=(\alph*)]
    \item 1
    \item 2
    \item 3
    \item 5
\end{enumerate}
\item For which of the following thermodynamic processes is the increase in the internal energy of an ideal gas equal to the heat added to the gas?
\begin{enumerate}[label=(\alph*)]
    \item Constant temperature
    \item Constant volume
    \item Constant pressure
    \item Adiabatic
\end{enumerate}
\end{enumerate}

\section*{True / False (10 marks)}
\textit{Answer True or False}
\begin{enumerate}[label=\arabic*.]
\setcounter{enumi}{5}
\item True or False: According to BCS theory, Cooper pairs are attracted to each other through the strong nuclear force.
\item True or False: If a photon detector has a quantum efficiency of 0.1, it will detect an average of 10 photons with an rms deviation of about 3 from 100 sent photons.
\item True or False: An allowed transition from the n = 4, l = 1 state in hydrogen cannot reach the n = 3, l = 1 state.
\item True or False: The speed of light in a nonmagnetic dielectric material with a dielectric constant of 4.0 is 1.5 x $10^8$ m/s.
\item True or False: The primary source of the Sun's energy is the fusion of four hydrogen atoms into one helium atom, releasing energy according to E=$mc^2$.
\end{enumerate}

\section*{Long Form Response (20 marks)}
\textit{Answer each question with a clear and complete explanation.}
\begin{enumerate}[label=\arabic*.]
\setcounter{enumi}{10}
\item Discuss how Fermat's principle of ray optics, which states that light travels the path that requires the least time, can be used to derive both Snell's law of refraction and the law of reflection.
\item Discuss the work required to accelerate a particle of mass m from rest to a speed of 0.6 c, and analyze the implications of relativistic effects on this energy calculation.
\end{enumerate}

\vfill
\noindent\textit{End of Paper}
\end{document}